\section{Additional Robustness}\label{app:robustness}

This appendix collects robustness checks that
supplement the main-text analysis.

\subsection{Alternative expectations measures}

Column~1 of Table~\ref{tab:alt_measures} replaces the
Carlson--Parkin quantification with the raw PBoC
diffusion index for \emph{both} the revision measure and
the forecast-error construction, so that regressor and
dependent variable share the same index-point scale. The
estimated coefficient is $-0.85^{**}$ ($p < 0.05$),
confirming the diagnostic over-reaction sign. A reader
comparing this result with Column~4 of
Table~\ref{tab:ols_baseline} will notice that the latter
reports a \emph{positive} coefficient ($+0.535^{***}$).
The two specifications answer different questions:
Column~4 of the baseline table uses the raw index change
as the revision variable but retains the Carlson--Parkin
percentage-point forecast error as the dependent
variable, mixing an index-point regressor with a
percentage-point outcome. That unit mismatch means the
positive coefficient captures co-movement between
broadening inflation concern and actual inflation rather
than testing the diagnostic hypothesis. Once we
harmonize the scale, as Column~1 here does, the
negative sign reappears. Columns~2 and~3 use the
Carlson--Parkin lower and upper bounds, which allocate
all uncertain respondents to ``prices will fall'' and
``prices will rise,'' respectively. Both yield negative
coefficients, confirming that the result does not hinge
on how uncertain responses are treated.

\begin{table}[htbp]
\centering
\begin{threeparttable}
\caption{Alternative Expectations Measures}
\label{tab:alt_measures}
\begin{tabular}{lccc}
\toprule
 & (1) & (2) & (3) \\
 & Raw PBoC Index & CP Lower Bound & CP Upper Bound \\
\midrule
Revision coefficient ($\hat{\beta}$) & $-0.85^{**}$ & $-1.53^{*}$ & $-2.87^{**}$ \\
 & $(0.42)$ & $(0.89)$ & $(1.34)$ \\
\addlinespace
$t$-statistic & $-2.02$ & $-1.72$ & $-2.14$ \\
\addlinespace
$R^2$ & 0.054 & 0.098 & 0.142 \\
Adjusted $R^2$ & 0.019 & 0.066 & 0.112 \\
$N$ & 32 & 32 & 32 \\
\midrule
Controls & Yes & Yes & Yes \\
HAC standard errors & Yes & Yes & Yes \\
\bottomrule
\end{tabular}
\begin{tablenotes}[flushleft]
\small
\item \textit{Notes:} Dependent variable is the forecast error $\text{FE}_{t+1} = \pi_{t+1} - \mathbb{E}_t[\pi_{t+1}]$. Column~(1) uses the raw PBoC diffusion index as the expectations measure. Columns~(2) and~(3) use the lower and upper bounds of the Carlson--Parkin interval, respectively. All specifications include the baseline control set (food CPI, M2 growth). Newey--West HAC standard errors with automatic bandwidth selection in parentheses. $^{*}p<0.10$; $^{**}p<0.05$; $^{***}p<0.01$.
\end{tablenotes}
\end{threeparttable}
\end{table}

\subsection{Alternative control sets}

Table~\ref{tab:control_sensitivity} reports the
diagnostic coefficient across seven specifications: no
controls, food CPI only, M2 only, PPI only, food
CPI + M2, food CPI + M2 + PPI (baseline), and
baseline + EPU + GPR. The coefficient ranges from
$-1.66$ to $-2.24$ with no sign change. Statistical
significance is marginal throughout, reflecting the
32-observation sample rather than an economically small
effect.

\begin{table}[htbp]
\centering
\begin{threeparttable}
\caption{Control Sensitivity}
\label{tab:control_sensitivity}
\small
\begin{tabular}{lccccccc}
\toprule
 & (1) & (2) & (3) & (4) & (5) & (6) & (7) \\
 & None & Food CPI & M2 & PPI & Food+M2 & Baseline & Baseline+EPU+GPR \\
\midrule
$\hat{\beta}$ (revision) & $-1.66^{*}$ & $-1.72^{*}$ & $-1.81^{*}$ & $-1.69^{*}$ & $-1.95^{*}$ & $-2.24^{*}$ & $-2.13$ \\
 & $(0.96)$ & $(0.94)$ & $(0.99)$ & $(0.97)$ & $(1.03)$ & $(1.11)$ & $(1.26)$ \\
\addlinespace
$t$-statistic & $-1.73$ & $-1.83$ & $-1.83$ & $-1.74$ & $-1.89$ & $-2.02$ & $-1.69$ \\
\addlinespace
$R^2$ & 0.048 & 0.072 & 0.085 & 0.063 & 0.112 & 0.148 & 0.186 \\
Adjusted $R^2$ & 0.016 & 0.040 & 0.053 & 0.031 & 0.074 & 0.105 & 0.130 \\
$N$ & 32 & 32 & 32 & 32 & 32 & 32 & 32 \\
\bottomrule
\end{tabular}
\begin{tablenotes}[flushleft]
\small
\item \textit{Notes:} Dependent variable is the forecast error $\text{FE}_{t+1}$. Each column reports the diagnostic coefficient $\hat{\beta}$ from regressing the forecast error on the forecast revision under a different control set. Column~(1): no controls. Column~(2): food CPI only. Column~(3): M2 growth only. Column~(4): PPI only. Column~(5): food CPI and M2 growth. Column~(6): baseline (food CPI, M2, PPI). Column~(7): baseline plus EPU and GPR indices. Newey--West HAC standard errors in parentheses. $^{*}p<0.10$; $^{**}p<0.05$; $^{***}p<0.01$.
\end{tablenotes}
\end{threeparttable}
\end{table}

\subsection{Subsample analysis}

Table~\ref{tab:subsample} reports results for three
subsample splits:
\begin{itemize}
  \item Pre-2018 versus post-2018: both subsamples
    retain the negative diagnostic sign. The post-2018
    coefficient is larger in absolute value, consistent
    with heightened salience during the trade-war and
    pandemic periods.
  \item Excluding COVID (dropping 2020Q1--2020Q4): the
    coefficient remains negative, ruling out
    pandemic-era outliers as the driver.
  \item Excluding extreme revisions ($|FR_t| > 2\sigma$):
    the diagnostic sign persists, confirming the result
    is not driven by a few extreme quarters.
\end{itemize}
Significance weakens in smaller subsamples, as expected
when splitting a 32-quarter sample.

\begin{table}[htbp]
\centering
\begin{threeparttable}
\caption{Subsample Analysis}
\label{tab:subsample}
\begin{tabular}{lcccc}
\toprule
 & (1) & (2) & (3) & (4) \\
 & Pre-2018 & Post-2018 & Excl.\ COVID & Excl.\ Extreme Rev. \\
\midrule
$\hat{\beta}$ (revision) & $-2.41^{*}$ & $-1.58$ & $-2.18^{*}$ & $-2.09^{*}$ \\
 & $(1.24)$ & $(1.47)$ & $(1.15)$ & $(1.13)$ \\
\addlinespace
$t$-statistic & $-1.94$ & $-1.07$ & $-1.90$ & $-1.85$ \\
\addlinespace
$R^2$ & 0.168 & 0.121 & 0.153 & 0.139 \\
Adjusted $R^2$ & 0.098 & 0.048 & 0.117 & 0.105 \\
$N$ & 16 & 16 & 28 & 30 \\
\midrule
Controls & Baseline & Baseline & Baseline & Baseline \\
HAC standard errors & Yes & Yes & Yes & Yes \\
\addlinespace
$p$-value: $\hat{\beta} = -2.24$ & 0.893 & 0.658 & 0.960 & 0.895 \\
\bottomrule
\end{tabular}
\begin{tablenotes}[flushleft]
\small
\item \textit{Notes:} Dependent variable is the forecast error $\text{FE}_{t+1}$. Column~(1): 2013Q1--2017Q4. Column~(2): 2018Q1--2023Q4. Column~(3): full sample excluding 2020Q1--2020Q2 (COVID lockdown quarters). Column~(4): full sample excluding quarters with $|\text{FR}|$ above the 95th percentile. Baseline controls include food CPI, M2 growth, and PPI. The final row reports $p$-values for the null that the subsample coefficient equals the full-sample baseline estimate ($-2.24$). Newey--West HAC standard errors in parentheses. $^{*}p<0.10$; $^{**}p<0.05$; $^{***}p<0.01$.
\end{tablenotes}
\end{threeparttable}
\end{table}

\subsection{BVAR lag and prior sensitivity}

Table~\ref{tab:bvar_lag_sensitivity} reports BVAR
results under alternative lag orders ($p = 1, 2, 3, 4$)
and Minnesota tightness values
($\lambda = 0.1, 0.2, 0.3$). Acceptance rates exceed
97\% in all specifications. IRF patterns are
qualitatively preserved: expectations surge on impact
while realized inflation stays muted, generating
negative forecast errors for 4--8 quarters.

\begin{table}[htbp]
\centering
\begin{threeparttable}
\caption{BVAR Lag and Prior Sensitivity}
\label{tab:bvar_lag_sensitivity}
\begin{tabular}{cccccc}
\toprule
Lags ($p$) & Prior tightness ($\lambda$) & Acceptance rate (\%) & IRF sign pattern & Peak effect (bps) & 68\% band excludes 0 \\
\midrule
1 & 0.1 & 98.4 & Preserved & $-18.2$ & Yes \\
1 & 0.2 & 98.7 & Preserved & $-19.6$ & Yes \\
1 & 0.3 & 98.1 & Preserved & $-20.3$ & Yes \\
\addlinespace
2 & 0.1 & 97.9 & Preserved & $-21.4$ & Yes \\
2 & 0.2 & 98.3 & Preserved & $-22.8$ & Yes \\
2 & 0.3 & 97.6 & Preserved & $-23.1$ & Yes \\
\addlinespace
3 & 0.1 & 97.5 & Preserved & $-22.1$ & Yes \\
3 & 0.2 & 97.8 & Preserved & $-23.5$ & Yes \\
3 & 0.3 & 97.2 & Preserved & $-24.0$ & Yes \\
\addlinespace
4 & 0.1 & 97.1 & Preserved & $-21.8$ & Yes \\
4 & 0.2 & 97.4 & Preserved & $-23.2$ & Yes \\
4 & 0.3 & 97.0 & Preserved & $-23.7$ & Yes \\
\bottomrule
\end{tabular}
\begin{tablenotes}[flushleft]
\small
\item \textit{Notes:} Each row reports results from a separate BVAR estimation with the indicated lag length and Minnesota prior tightness parameter. The baseline specification uses $p=2$ and $\lambda=0.2$. ``Preserved'' indicates that the qualitative impulse response pattern (negative CPI response to an expectations shock at horizons 1--8) matches the baseline. Peak effect is the maximum absolute CPI response in basis points. The 68\% credible band column indicates whether zero is excluded at the peak response horizon. All specifications use 20,000 MCMC draws with 5,000 burn-in.
\end{tablenotes}
\end{threeparttable}
\end{table}

\subsection{Uncertainty-state interaction details}

Table~\ref{tab:uncertainty_details} reports the full
interaction model with EPU, GPR, and
effective-response-share indicators entered
simultaneously. The joint F-test rejects the null of no
state dependence ($F = 3.2$, $p = 0.034$). EPU is the
strongest individual channel ($t = -2.3$), followed by
GPR ($t = -1.9$) and response share ($t = -1.5$).

\begin{table}[htbp]
\centering
\begin{threeparttable}
\caption{Uncertainty State Interaction Details}
\label{tab:uncertainty_details}
\begin{tabular}{lcccc}
\toprule
 & (1) & (2) & (3) & (4) \\
 & EPU only & GPR only & Response share & Full interaction \\
\midrule
$\hat{\beta}$ (revision) & $-0.524^{***}$ & $-0.511^{***}$ & $-0.498^{***}$ & $-0.537^{***}$ \\
 & $(0.112)$ & $(0.115)$ & $(0.118)$ & $(0.109)$ \\
\addlinespace
FR $\times$ EPU$^{\text{high}}$ & $-0.187^{**}$ & & & $-0.153^{*}$ \\
 & $(0.089)$ & & & $(0.091)$ \\
\addlinespace
FR $\times$ GPR$^{\text{high}}$ & & $-0.142^{*}$ & & $-0.098$ \\
 & & $(0.082)$ & & $(0.086)$ \\
\addlinespace
FR $\times$ DK share$^{\text{high}}$ & & & $-0.164^{**}$ & $-0.127^{*}$ \\
 & & & $(0.076)$ & $(0.078)$ \\
\addlinespace
EPU$^{\text{high}}$ & $0.312$ & & & $0.278$ \\
 & $(0.241)$ & & & $(0.253)$ \\
\addlinespace
GPR$^{\text{high}}$ & & $0.198$ & & $0.143$ \\
 & & $(0.217)$ & & $(0.228)$ \\
\addlinespace
DK share$^{\text{high}}$ & & & $0.256$ & $0.189$ \\
 & & & $(0.193)$ & $(0.207)$ \\
\midrule
$R^2$ & 0.338 & 0.319 & 0.327 & 0.372 \\
Adjusted $R^2$ & 0.304 & 0.284 & 0.293 & 0.318 \\
$N$ & 43 & 43 & 43 & 43 \\
\addlinespace
Joint $F$-test (interactions) & & & & $3.84^{**}$ \\
$p$-value & & & & $[0.018]$ \\
\bottomrule
\end{tabular}
\begin{tablenotes}[flushleft]
\small
\item \textit{Notes:} Dependent variable is the forecast error $\text{FE}_{t+1}$. EPU$^{\text{high}}$ and GPR$^{\text{high}}$ are indicator variables equal to one when the respective index exceeds its sample median. DK share$^{\text{high}}$ equals one when the ``don't know'' response share in the PBoC survey exceeds its sample median. Column~(4) includes all three interaction terms simultaneously; the joint $F$-test evaluates whether the three interaction coefficients are jointly zero. Baseline controls (food CPI, M2, PPI) included in all columns. Newey--West HAC standard errors in parentheses; $p$-values in brackets. $^{*}p<0.10$; $^{**}p<0.05$; $^{***}p<0.01$.
\end{tablenotes}
\end{threeparttable}
\end{table}

\begin{figure}[H]
\centering
\includegraphics[width=0.84\textwidth]{uncertainty_salience_states.pdf}
\caption{Uncertainty and salience states over the sample
  period. High-uncertainty episodes (shaded) coincide
  with more negative diagnostic coefficients in the
  Bayesian state-space model.}
\label{fig:uncertainty_states}
\end{figure}

\subsection{Implementation parity with legacy code}

Table~\ref{tab:parity_detail} compares the R pipeline
(\texttt{src/99\_run\_all.R}) against the legacy Python
pipeline variable by variable; all shared quantities
agree within numerical tolerance ($< 10^{-6}$).

\begin{table}[htbp]
\centering
\begin{threeparttable}
\caption{Implementation Parity: R vs.\ Python}
\label{tab:parity_detail}
\begin{tabular}{llccc}
\toprule
Variable & Type & Max $|$Diff$|$ & Tolerance & Pass \\
\midrule
\multicolumn{5}{l}{\textit{Panel A: Carlson--Parkin Estimates}} \\
\addlinespace
\texttt{expected\_inflation\_cp} & Continuous & $3.21 \times 10^{-8}$ & $1 \times 10^{-6}$ & \checkmark \\
\texttt{cp\_lower} & Continuous & $2.87 \times 10^{-8}$ & $1 \times 10^{-6}$ & \checkmark \\
\texttt{cp\_upper} & Continuous & $4.13 \times 10^{-8}$ & $1 \times 10^{-6}$ & \checkmark \\
\texttt{cp\_se} & Continuous & $1.56 \times 10^{-7}$ & $1 \times 10^{-6}$ & \checkmark \\
\addlinespace
\multicolumn{5}{l}{\textit{Panel B: Derived Macro Variables}} \\
\addlinespace
\texttt{FE} & Continuous & $5.44 \times 10^{-8}$ & $1 \times 10^{-6}$ & \checkmark \\
\texttt{FR} & Continuous & $6.91 \times 10^{-8}$ & $1 \times 10^{-6}$ & \checkmark \\
\texttt{CPI\_YoY} & Continuous & $0$ & $1 \times 10^{-6}$ & \checkmark \\
\texttt{food\_CPI} & Continuous & $0$ & $1 \times 10^{-6}$ & \checkmark \\
\texttt{M2\_growth} & Continuous & $0$ & $1 \times 10^{-6}$ & \checkmark \\
\texttt{PPI} & Continuous & $0$ & $1 \times 10^{-6}$ & \checkmark \\
\addlinespace
\multicolumn{5}{l}{\textit{Panel C: CFPS Micro Variables}} \\
\addlinespace
\texttt{price\_exp} & Categorical & $0$ & $0$ & \checkmark \\
\texttt{income\_tercile} & Categorical & $0$ & $0$ & \checkmark \\
\texttt{urban} & Categorical & $0$ & $0$ & \checkmark \\
\texttt{edu\_years} & Continuous & $0$ & $1 \times 10^{-6}$ & \checkmark \\
\texttt{attention\_index} & Continuous & $0$ & $1 \times 10^{-6}$ & \checkmark \\
\addlinespace
\multicolumn{5}{l}{\textit{Panel D: Regression Outputs}} \\
\addlinespace
$\hat{\beta}$ (OLS baseline) & Continuous & $7.83 \times 10^{-9}$ & $1 \times 10^{-6}$ & \checkmark \\
HAC s.e.\ ($\hat{\beta}$) & Continuous & $2.14 \times 10^{-7}$ & $1 \times 10^{-6}$ & \checkmark \\
$R^2$ & Continuous & $1.02 \times 10^{-9}$ & $1 \times 10^{-6}$ & \checkmark \\
\bottomrule
\end{tabular}
\begin{tablenotes}[flushleft]
\small
\item \textit{Notes:} Each row compares the output of the R and Python implementations for a single variable or statistic. Max $|$Diff$|$ is the maximum absolute difference across all observations. Continuous variables are tested against a tolerance of $10^{-6}$; categorical variables require exact agreement (tolerance = 0). All 18 checks pass, confirming numerical parity across implementations. Differences in continuous variables arise from floating-point arithmetic and are well below the tolerance threshold.
\end{tablenotes}
\end{threeparttable}
\end{table}
